\documentclass[12pt,a4paper]{report}
\usepackage[utf8]{inputenc}
\usepackage[spanish]{babel}
\usepackage{amsmath, amssymb}
\usepackage{graphicx}
\usepackage{booktabs}
\usepackage{geometry}
\usepackage{hyperref}
\usepackage{tikz}
\usepackage{float}
\usepackage{fancyhdr}
\usepackage{titlesec}
\usepackage{setspace}

\geometry{left=3cm, right=2.5cm, top=3cm, bottom=3cm}

\hypersetup{
    colorlinks=true,
    linkcolor=blue,
    filecolor=magenta,      
    urlcolor=cyan,
    pdftitle={Informe: Sistema Híbrido de Recomendación},
}

\begin{document}

% -------------------------------------------------------------------
% PORTADA
% -------------------------------------------------------------------
\begin{titlepage}
    \centering
    {\bfseries\LARGE UNIVERSIDAD NACIONAL DEL ALTIPLANO \par}
    \vspace{0.5cm}
    {\scshape\Large Facultad de Ingeniería Mecánica Eléctrica, Electrónica y Sistemas \par}
    \vspace{0.3cm}
    {\scshape\large Escuela Profesional de Ingeniería de Sistemas \par}
    
    \vspace{1.5cm}
    % Placeholder para el logo. Puedes reemplazar el comando \rule por \includegraphics{logo.png} cuando tengas tu imagen
    \begin{tikzpicture}
        \draw[thick, rounded corners, fill=gray!10] (0,0) rectangle (4,4);
        \node at (2,2) {\textbf{[LOGO UNAP AQUÍ]}};
        \node at (2,1.5) {\footnotesize (Reemplazar en el .tex)};
    \end{tikzpicture}
    % \includegraphics[width=0.3\textwidth]{logo.png}
    
    \vspace{1.5cm}
    {\scshape\Large Informe del Proyecto Final \par}
    \vspace{0.5cm}
    {\huge\bfseries Sistema Híbrido de Recomendación de Laptops en Dos Niveles para Mitigación de \textit{Cold Start} \par}
    
    \vspace{2cm}
    \begin{flushleft}
        \large
        \textbf{Curso:} Sistemas de Recomendación \\
        \textbf{Docente:} Mgtr. GONZALES PACO MAGALI GIANINA \\
        \vspace{0.5cm}
        \textbf{Integrantes:}\\
        \indent $\bullet$ Apaza Condori Carmen Nieves \\
        \indent $\bullet$ Vilca Cari Aaron Roger
    \end{flushleft}
    
    \vfill
    {\large Puno - Perú \par}
    {\large 2026 \par}
\end{titlepage}

% -------------------------------------------------------------------
% ÍNDICE Y FORMATO
% -------------------------------------------------------------------
\setstretch{1.15}
\tableofcontents
\newpage

\pagestyle{fancy}
\fancyhf{}
\rhead{Sistema Híbrido de Recomendación}
\lhead{UNAP - Ingeniería de Sistemas}
\cfoot{\thepage}

% -------------------------------------------------------------------
% CAPÍTULO 1: INTRODUCCIÓN
% -------------------------------------------------------------------
\chapter{Introducción}
\section{Contexto del Problema}
En la era del comercio electrónico (E-Commerce), los sistemas de recomendación (SR) son fundamentales para guiar a los consumidores frente a la sobreabundancia de opciones. Plataformas como Amazon o Netflix dependen en gran medida del Filtrado Colaborativo. Sin embargo, cuando se trata de hardware tecnológico (como las computadoras portátiles o laptops), la naturaleza del dominio presenta desafíos únicos que hacen que los modelos tradicionales fracasen.

El principal problema es el \textbf{Inicio en Frío (Cold Start)}. Las computadoras no son artículos de consumo diario; un usuario promedio adquiere una nueva laptop cada 3 o 5 años. Por ende, la inmensa mayoría de las sesiones en un sistema de ventas de hardware carecen de historial de calificaciones previo. Si un algoritmo se basa únicamente en colaboraciones (compras de otros usuarios similares), fallará estrepitosamente al recomendar equipos a un cliente completamente nuevo.

A esto se suma la complejidad técnica de las restricciones duras (\textit{Hard Constraints}). A diferencia de una canción o película, donde un falso positivo solo genera aburrimiento, en hardware un equipo incompatible (por ejemplo, recomendar una laptop sin tarjeta de video dedicada a un diseñador 3D) invalida completamente la transacción y destruye la confianza del cliente.

\section{Justificación}
Es imperativo desarrollar una arquitectura inteligente que no solo entienda los gustos subjetivos, sino que evalúe lógicamente los requisitos de hardware. Este proyecto aborda dichos desafíos diseñando un **Sistema Híbrido de Recomendación de Dos Niveles**, el cual prioriza las matemáticas de la Teoría de Utilidad Multi-Atributo (MAUT) y la Lógica Condicional frente al modelo colaborativo clásico, logrando sortear el vacío de datos iniciales.

% -------------------------------------------------------------------
% CAPÍTULO 2: OBJETIVOS
% -------------------------------------------------------------------
\chapter{Objetivos}
\section{Objetivo General}
Diseñar, implementar y evaluar un Sistema Híbrido de Recomendación de Laptops estructurado en dos niveles computacionales para maximizar la tasa de satisfacción de restricciones y mitigar el problema de \textit{Cold Start} en entornos de E-Commerce tecnológico.

\section{Objetivos Específicos}
\begin{itemize}
    \item Procesar e integrar un dataset real de especificaciones de laptops (proveniente de Kaggle) extrayendo características vitales mediante expresiones regulares (RAM, CPU, GPU).
    \item Diseñar un Motor de Inferencia Lógica (Nivel 1) para aplicar filtros duros que descarten hardware incompetente frente a un perfil de uso dado.
    \item Desarrollar un ensamble dinámico ponderado (Nivel 2) compuesto por MAUT, Factorización de Matrices (SVD) y Similitud Coseno.
    \item Construir un Dashboard en un entorno web que permita la interacción en tiempo real del modelo y visualice la adaptación de pesos $\alpha, \beta, \gamma$.
    \item Evaluar el rendimiento del sistema contrastando su métrica de Precisión y \textit{Constraint Satisfaction Rate} (CSR) frente a un Baseline de SVD Clásico.
\end{itemize}

% -------------------------------------------------------------------
% CAPÍTULO 3: MARCO TEÓRICO
% -------------------------------------------------------------------
\chapter{Marco Teórico}
\section{Sistemas de Recomendación Colaborativos (SVD)}
El filtrado colaborativo asume que si dos usuarios compartieron intereses en el pasado, lo harán en el futuro. La técnica más prominente es la \textbf{Factorización de Matrices}, utilizando \textit{Singular Value Decomposition (SVD)}. Este enfoque estima una calificación $\hat{r}_{ui}$ como el producto escalar de dos vectores de factores latentes (usuario $p_u$ e ítem $q_i$), ajustado por sesgos (\textit{bias}):
\begin{equation}
    \hat{r}_{ui} = \mu + b_u + b_i + q_i^T p_u
\end{equation}
Este modelo se optimiza usualmente mediante Descenso de Gradiente Estocástico (SGD). A pesar de su efectividad geométrica, es inútil frente al fenómeno de \textit{Cold Start}.

\section{Recomendadores Basados en Conocimiento y MAUT}
Los sistemas Knowledge-Based no dependen del historial, sino del emparejamiento directo entre los requisitos declarados por el usuario y la ontología del producto. La \textbf{Teoría de la Utilidad Multi-Atributo (MAUT)} es un enfoque matemático que permite puntuar un ítem evaluando ponderadamente cada una de sus características normalizadas. Para una laptop, los atributos ($S$) como RAM, Procesador o Precio se estandarizan a valores entre 0 y 1, multiplicándose por la preferencia del usuario ($w$).

\section{Recomendadores Basados en Contenido}
El enfoque Content-Based recomienda ítems similares a los preferidos anteriormente. Matemáticamente se calcula construyendo un vector de características para la laptop $v_i$ y comparándolo con el vector ideal del usuario $v_u$ usando \textbf{Similitud Coseno}:
\begin{equation}
    \text{Coseno}(v_u, v_i) = \frac{v_u \cdot v_i}{\|v_u\| \|v_i\|}
\end{equation}

% -------------------------------------------------------------------
% CAPÍTULO 4: ARQUITECTURA DEL SISTEMA
% -------------------------------------------------------------------
\chapter{Arquitectura del Sistema Híbrido}
El núcleo de la investigación reside en la arquitectura en cascada de dos niveles, ideada para bloquear inmediatamente anomalías operativas.

\section{Nivel 1: Inferencia Lógica (Constraint-Based Filtering)}
Actúa como un firewall técnico. Cuando un usuario introduce su presupuesto y selecciona un "Perfil de Uso" (ej. Arquitectura 3D), el motor extrae de su base de conocimiento las reglas mínimas.
\begin{equation}
    M_{rules}(i) = \begin{cases} 
    1 & \text{si } P_i \le B \land (R_i \ge R_{min}) \land (V_i \ge V_{min}) \\
    0 & \text{en cualquier otro caso}
    \end{cases}
\end{equation}
Cualquier equipo con $M_{rules}=0$ se anula para la segunda fase.

\section{Nivel 2: Ensamble Dinámico Ponderado}
Aplica tres sub-modelos a los equipos sobrevivientes y unifica sus resultados usando ponderaciones adaptativas:
\begin{equation}
    Score_i = M_{rules}(i) \times \Big[ \alpha \cdot U_{maut}(i) + \beta \cdot \hat{r}_{svd}(i) + \gamma \cdot S_{content}(i) \Big]
\end{equation}
Donde $\alpha + \beta + \gamma = 1$. 

\textbf{Gestión de Pesos:}
Para mitigar el \textit{Cold Start}, el sistema altera los pesos dinámicamente:
\begin{itemize}
    \item \textit{Usuario Nuevo}: Se configura $\beta = 0.10$ (se ignora el SVD) y se maximiza $\alpha = 0.60$ para usar MAUT.
    \item \textit{Usuario Consolidado}: Se balancean los pesos para aprovechar el filtrado colaborativo de la comunidad ($\beta = 0.50$).
\end{itemize}

% -------------------------------------------------------------------
% CAPÍTULO 5: METODOLOGÍA
% -------------------------------------------------------------------
\chapter{Desarrollo y Metodología}
\section{Dataset y Preprocesamiento}
Se usó un dataset de Kaggle (2024) con 1,273 modelos reales de laptops.
El procesamiento (realizado en Python con Pandas) incluyó:
\begin{itemize}
    \item Limpieza de cadenas de texto y extracción (Regex) de valores como `16GB RAM` o `8 Cores`.
    \item Detección de tarjetas dedicadas (Nvidia, AMD Radeon) para activar el flag booleano de soporte CUDA y render.
    \item Conversión de moneda Rupias Indias (INR) a Dólares (USD).
\end{itemize}

\section{Implementación API y Dashboard}
La lógica algorítmica fue empaquetada en una API REST construida con \textbf{FastAPI}. Esta API alimenta a un Dashboard de usuario implementado con HTML5, CSS3 y JavaScript asíncrono.
El frontend permite manipular deslizadores (sliders) para configurar los pesos MAUT y permite observar en vivo cómo los gráficos estadísticos y las recomendaciones varían en milisegundos.

% -------------------------------------------------------------------
% CAPÍTULO 6: RESULTADOS
% -------------------------------------------------------------------
\chapter{Evaluación y Resultados}
Para demostrar la eficacia técnica del ensamblaje híbrido, se calcularon métricas contra dos baselines: un modelo colaborativo (SVD Puro) y un modelo basado en contenido puro.

\section{Constraint Satisfaction Rate (CSR@10)}
Esta métrica, vital en el comercio de hardware, evalúa qué porcentaje de los equipos presentados en el Top-10 cumplen con las especificaciones mínimas requeridas.
\begin{table}[H]
    \centering
    \begin{tabular}{lc}
        \toprule
        \textbf{Modelo Analizado} & \textbf{CSR@10} \\
        \midrule
        Baseline SVD (Colaborativo) & 31.7\% \\
        Baseline Similitud (Content-Based) & 48.3\% \\
        \textbf{Sistema Híbrido Cascada Propuesto} & \textbf{93.3\%} \\
        \bottomrule
    \end{tabular}
    \caption{Resultados Comparativos de Satisfacción de Restricciones}
\end{table}
El modelo híbrido superó contundentemente a la IA tradicional. El SVD clásico falló estrepitosamente debido al "sesgo de popularidad", recomendando equipos de ofimática a diseñadores solo porque esos equipos se venden masivamente.

\section{Error Cuadrático Medio (RMSE)}
El componente SVD interno del ensamble híbrido se testeó independientemente en un split 80/20, obteniendo un RMSE de \textbf{1.033}. Esto confirma que el modelo latente logró aprender a estimar calificaciones con una desviación de 1 estrella.

% -------------------------------------------------------------------
% CAPÍTULO 7: CONCLUSIONES
% -------------------------------------------------------------------
\chapter{Conclusiones}
\begin{itemize}
    \item \textbf{Superación de Limitaciones Tradicionales:} Se comprobó empíricamente que el uso aislado del Filtrado Colaborativo (como el SVD) no es apto para plataformas de ventas técnicas donde las restricciones de hardware determinan la usabilidad.
    \item \textbf{Éxito en el Cold Start:} La incorporación de la Teoría MAUT demostró ser un puente matemático infalible para valorar ítems a través de atributos puros en ausencia de historial de consumidor, reflejado en un aumento del CSR al 93.3\%.
    \item \textbf{Eficiencia Computacional:} La inclusión de una Máscara Binaria (Nivel 1) no solo aporta seguridad lógica, sino que disminuye drásticamente el espacio de cálculo para la similitud coseno y la factorización en el Nivel 2.
\end{itemize}

% -------------------------------------------------------------------
% REFERENCIAS
% -------------------------------------------------------------------
\begin{thebibliography}{9}
\bibitem{ricci} 
Ricci, F., Rokach, L., \& Shapira, B. (2015). 
\textit{Recommender Systems Handbook}. Springer.

\bibitem{koren} 
Koren, Y., Bell, R., \& Volinsky, C. (2009). 
Matrix Factorization Techniques for Recommender Systems. \textit{Computer, 42}(8), 30-37.

\bibitem{burke} 
Burke, R. (2002). 
Hybrid Recommender Systems: Survey and Experiments. \textit{User Modeling and User-Adapted Interaction, 12}(4), 331-370.
\end{thebibliography}

\end{document}
